\documentclass[10pt,a4paper]{amsart}
\usepackage{amsmath,amssymb,amsthm,graphicx}
\usepackage{verbatim}
\newenvironment{reference}{\comment}{\endcomment}
%\newenvironment{reference}{}{}
\newenvironment{slogan}{\comment}{\endcomment}
\newenvironment{history}{\comment}{\endcomment}

% For commutative diagrams we use Xy-pic
\usepackage[all]{xy}

% We use 2cell for 2-commutative diagrams.
\xyoption{2cell}
% Theorem environments.
%
\theoremstyle{plain}
\newtheorem{theorem}[subsection]{Theorem}
\newtheorem{proposition}[subsection]{Proposition}
\newtheorem{lemma}[subsection]{Lemma}

\theoremstyle{definition}
\newtheorem{definition}[subsection]{Definition}
\newtheorem{example}[subsection]{Example}
\newtheorem{exercise}[subsection]{Exercise}
\newtheorem{situation}[subsection]{Situation}

\theoremstyle{remark}
\newtheorem{remark}[subsection]{Remark}
\newtheorem{remarks}[subsection]{Remarks}

\numberwithin{equation}{subsection}

% Macros
%
\def\lim{\mathop{\mathrm{lim}}\nolimits}
\def\colim{\mathop{\mathrm{colim}}\nolimits}
\def\Spec{\mathop{\mathrm{Spec}}}
\def\Hom{\mathop{\mathrm{Hom}}\nolimits}
\def\Ext{\mathop{\mathrm{Ext}}\nolimits}
\def\SheafHom{\mathop{\mathcal{H}\!\mathit{om}}\nolimits}
\def\SheafExt{\mathop{\mathcal{E}\!\mathit{xt}}\nolimits}
\def\Sch{\mathit{Sch}}
\def\Mor{\mathop{\mathrm{Mor}}\nolimits}
\def\Ob{\mathop{\mathrm{Ob}}\nolimits}
\def\Sh{\mathop{\mathit{Sh}}\nolimits}
\def\NL{\mathop{N\!L}\nolimits}
\def\CH{\mathop{\mathrm{CH}}\nolimits}
\def\proetale{{pro\text{-}\acute{e}tale}}
\def\etale{{\acute{e}tale}}
\def\QCoh{\mathit{QCoh}}
\def\Ker{\mathop{\mathrm{Ker}}}
\def\Im{\mathop{\mathrm{Im}}}
\def\Coker{\mathop{\mathrm{Coker}}}
\def\Coim{\mathop{\mathrm{Coim}}}

% Boxtimes
%
\DeclareMathSymbol{\boxtimes}{\mathbin}{AMSa}{"02}

%
% Macros for moduli stacks/spaces
%
\def\QCohstack{\mathcal{QC}\!\mathit{oh}}
\def\Cohstack{\mathcal{C}\!\mathit{oh}}
\def\Spacesstack{\mathcal{S}\!\mathit{paces}}
\def\Quotfunctor{\mathrm{Quot}}
\def\Hilbfunctor{\mathrm{Hilb}}
\def\Curvesstack{\mathcal{C}\!\mathit{urves}}
\def\Polarizedstack{\mathcal{P}\!\mathit{olarized}}
\def\Complexesstack{\mathcal{C}\!\mathit{omplexes}}
% \Pic is the operator that assigns to X its picard group, usage \Pic(X)
% \Picardstack_{X/B} denotes the Picard stack of X over B
% \Picardfunctor_{X/B} denotes the Picard functor of X over B
\def\Pic{\mathop{\mathrm{Pic}}\nolimits}
\def\Picardstack{\mathcal{P}\!\mathit{ic}}
\def\Picardfunctor{\mathrm{Pic}}
\def\Deformationcategory{\mathcal{D}\!\mathit{ef}}
\usepackage{cleveref}
\crefname{originalreference}{original reference}{original references}
\Crefname{originalreference}{Original reference}{Original references}
\makeatletter
\newlabel{remark-relative-dualizing-complex}{{0B6S}{source}{}{sourcepointer-1}{}}
\newlabel{remark-relative-dualizing-complex@cref}{{[originalreference][0][]0B6S}{[1][0][]source}{}{}{}}
\makeatother
\crefname{originallemma}{original lemma}{original lemmas}
\Crefname{originallemma}{Original lemma}{Original lemmas}
\makeatletter
\newlabel{cohomology-lemma-dual-perfect-complex}{{08DQ}{source}{}{sourcepointer-2}{}}
\newlabel{cohomology-lemma-dual-perfect-complex@cref}{{[originallemma][0][]08DQ}{[1][0][]source}{}{}{}}
\makeatother
\crefname{originallemma}{original lemma}{original lemmas}
\Crefname{originallemma}{Original lemma}{Original lemmas}
\makeatletter
\newlabel{lemma-iso-on-RSheafHom}{{0A9Q}{source}{}{sourcepointer-3}{}}
\newlabel{lemma-iso-on-RSheafHom@cref}{{[originallemma][0][]0A9Q}{[1][0][]source}{}{}{}}
\makeatother
\crefname{originalreference}{original reference}{original references}
\Crefname{originalreference}{Original reference}{Original references}
\makeatletter
\newlabel{remark-iso-on-RSheafHom}{{0GEV}{source}{}{sourcepointer-4}{}}
\newlabel{remark-iso-on-RSheafHom@cref}{{[originalreference][0][]0GEV}{[1][0][]source}{}{}{}}
\makeatother
\crefname{originallemma}{original lemma}{original lemmas}
\Crefname{originallemma}{Original lemma}{Original lemmas}
\makeatletter
\newlabel{perfect-lemma-flat-proper-perfect-direct-image-general}{{0B91}{source}{}{sourcepointer-5}{}}
\newlabel{perfect-lemma-flat-proper-perfect-direct-image-general@cref}{{[originallemma][0][]0B91}{[1][0][]source}{}{}{}}
\makeatother
\crefname{originallemma}{original lemma}{original lemmas}
\Crefname{originallemma}{Original lemma}{Original lemmas}
\makeatletter
\newlabel{more-morphisms-lemma-characterize-relatively-perfect}{{0GET}{source}{}{sourcepointer-6}{}}
\newlabel{more-morphisms-lemma-characterize-relatively-perfect@cref}{{[originallemma][0][]0GET}{[1][0][]source}{}{}{}}
\makeatother
\crefname{originallemma}{original lemma}{original lemmas}
\Crefname{originallemma}{Original lemma}{Original lemmas}
\makeatletter
\newlabel{cohomology-lemma-adjoint}{{079W}{source}{}{sourcepointer-7}{}}
\newlabel{cohomology-lemma-adjoint@cref}{{[originallemma][0][]079W}{[1][0][]source}{}{}{}}
\makeatother
\crefname{originallemma}{original lemma}{original lemmas}
\Crefname{originallemma}{Original lemma}{Original lemmas}
\makeatletter
\newlabel{perfect-lemma-cohomology-base-change}{{08EU}{source}{}{sourcepointer-8}{}}
\newlabel{perfect-lemma-cohomology-base-change@cref}{{[originallemma][0][]08EU}{[1][0][]source}{}{}{}}
\makeatother
\crefname{originallemma}{original lemma}{original lemmas}
\Crefname{originallemma}{Original lemma}{Original lemmas}
\makeatletter
\newlabel{more-algebra-lemma-perfect}{{0658}{source}{}{sourcepointer-9}{}}
\newlabel{more-algebra-lemma-perfect@cref}{{[originallemma][0][]0658}{[1][0][]source}{}{}{}}
\makeatother
\crefname{originallemma}{original lemma}{original lemmas}
\Crefname{originallemma}{Original lemma}{Original lemmas}
\makeatletter
\newlabel{cohomology-lemma-internal-hom}{{08DJ}{source}{}{sourcepointer-10}{}}
\newlabel{cohomology-lemma-internal-hom@cref}{{[originallemma][0][]08DJ}{[1][0][]source}{}{}{}}
\makeatother
\crefname{originallemma}{original lemma}{original lemmas}
\Crefname{originallemma}{Original lemma}{Original lemmas}
\makeatletter
\newlabel{perfect-lemma-quasi-coherence-internal-hom}{{0A6H}{source}{}{sourcepointer-11}{}}
\newlabel{perfect-lemma-quasi-coherence-internal-hom@cref}{{[originallemma][0][]0A6H}{[1][0][]source}{}{}{}}
\makeatother
\crefname{originaltheorem}{original theorem}{original theorems}
\Crefname{originaltheorem}{Original theorem}{Original theorems}
\makeatletter
\newlabel{perfect-theorem-bondal-van-den-Bergh}{{09IS}{source}{}{sourcepointer-12}{}}
\newlabel{perfect-theorem-bondal-van-den-Bergh@cref}{{[originaltheorem][0][]09IS}{[1][0][]source}{}{}{}}
\makeatother
\title{The relative dualizing complex of a proper flat finite presentation morphism is relatively perfect, has no positive cohomology after pushforward, and is self dual}
\author{The Stacks Project Authors}
\date{stacks-project:duality.tex:local-snapshot-20261012; GNU FDL 1.2 or later}
\begin{document}
\maketitle
\noindent\textit{Editorial note.} All delivered mathematical text is the author's own, copied from the pinned Stacks Project chapter file duality.tex; no statement or proof step was rewritten, reordered or removed.\par\medskip
\noindent\textit{Editorial note.} Preamble provenance: the pinned chapter file opens with an input of the shared preamble and carries no preamble text of its own, so the project's single Stacks-wide preamble staged at sources/stacks-duality/preamble.tex is excerpted verbatim under the recipe fixed by LANE\_BRIEF 22.1, as three spans, lines 8-13 (verbatim package and the reference/slogan/history environments), lines 14-18 (Xy-pic setup) and lines 154-217 (the native theorem declarations and mathematical macros). Nothing was added, removed or redefined; the wrapper class and the amsmath/amssymb/amsthm/graphicx base packages are supplied by the assembler.\par\medskip
\noindent\textit{Editorial note.} The counted statement and proof contain no cross reference, no citation and no included file, so the delivered body is self contained.\par\medskip
\noindent\textit{Editorial note.} Canonical location of the counted proposition: Stacks tag 0E4L (sources/stacks-duality/pinned-original-source.tex line 3028).\par\medskip
\noindent\textit{Allowed standard prerequisites (editorial).}\begin{itemize}
\item The relative dualizing complex and its base change formation, together with the right adjoint a of pushforward. Same chapter, Remark labelled remark-relative-dualizing-complex, cited at proof lines 3045-3049 and in the statement, and the construction of a used at proof lines 3076-3078. Supplies the object studied and allows the reduction to affine Y.
\item Duality for perfect complexes and the internal Hom. Cohomology, Lemmas labelled cohomology-lemma-dual-perfect-complex and cohomology-lemma-internal-hom, cited at proof lines 3055-3057, 3098-3100 and 3110-3112, and Cohomology, Lemma labelled cohomology-lemma-adjoint, cited at proof lines 3072-3074. Turns tensor products with the dualizing complex into duals and supplies the adjunction identities.
\item Perfectness of the pushforward of the structure sheaf and of the dualizing complex. Derived Categories of Schemes, Lemma labelled perfect-lemma-flat-proper-perfect-direct-image-general, cited at proof lines 3059-3063 and 3082-3086, and More on Morphisms, Lemma labelled more-morphisms-lemma-characterize-relatively-perfect, cited at proof lines 3068-3070. Gives the tor amplitude bound and converts the computation into Y-perfectness.
\item Generation by perfect objects and quasi-coherence of the internal Hom. Derived Categories of Schemes, Theorem labelled perfect-theorem-bondal-van-den-Bergh, cited at proof lines 3140-3144, and Lemma labelled perfect-lemma-quasi-coherence-internal-hom, cited at proof lines 3128-3130. Lets the functorial family of isomorphisms conclude that alpha is an isomorphism.
\item Representation of a perfect complex over an affine base. More on Algebra, Lemma labelled more-algebra-lemma-perfect, cited at proof lines 3088-3090. Represents Rf\_*O\_X by finite projective modules in degrees [0, N].
\end{itemize}\medskip
\ifdefined\hypertarget\hypertarget{sourcepointer-1}{}\fi
\noindent\textit{Source reference.} Original reference 0B6S: sources/stacks-duality/pinned-original-source.tex, source lines 2945--2945. Same-chapter pointer into the author environment remark at source line 2945; invoked at proof line 3047 of the counted proof. It is retained only as a navigation pointer; every load-bearing step is inside the delivered span.\par\medskip
\ifdefined\hypertarget\hypertarget{sourcepointer-2}{}\fi
\noindent\textit{Source reference.} Original lemma 08DQ: sources/stacks-chow/tags.txt, source lines 10645--10645. Cross-chapter standard prerequisite located in the project tag list staged in the local snapshot; invoked at proof line 3060 of the counted proof. It is retained only as a navigation pointer; every load-bearing step is inside the delivered span.\par\medskip
\ifdefined\hypertarget\hypertarget{sourcepointer-3}{}\fi
\noindent\textit{Source reference.} Original lemma 0A9Q: sources/stacks-duality/pinned-original-source.tex, source lines 775--775. Same-chapter pointer into the author environment lemma at source line 775; invoked at proof line 3061 of the counted proof. It is retained only as a navigation pointer; every load-bearing step is inside the delivered span.\par\medskip
\ifdefined\hypertarget\hypertarget{sourcepointer-4}{}\fi
\noindent\textit{Source reference.} Original reference 0GEV: sources/stacks-duality/pinned-original-source.tex, source lines 837--837. Same-chapter pointer into the author environment remark at source line 837; invoked at proof line 3062 of the counted proof. It is retained only as a navigation pointer; every load-bearing step is inside the delivered span.\par\medskip
\ifdefined\hypertarget\hypertarget{sourcepointer-5}{}\fi
\noindent\textit{Source reference.} Original lemma 0B91: sources/stacks-chow/tags.txt, source lines 14159--14159. Cross-chapter standard prerequisite located in the project tag list staged in the local snapshot; invoked at proof line 3064 of the counted proof. It is retained only as a navigation pointer; every load-bearing step is inside the delivered span.\par\medskip
\ifdefined\hypertarget\hypertarget{sourcepointer-6}{}\fi
\noindent\textit{Source reference.} Original lemma 0GET: sources/stacks-chow/tags.txt, source lines 20488--20488. Cross-chapter standard prerequisite located in the project tag list staged in the local snapshot; invoked at proof line 3069 of the counted proof. It is retained only as a navigation pointer; every load-bearing step is inside the delivered span.\par\medskip
\ifdefined\hypertarget\hypertarget{sourcepointer-7}{}\fi
\noindent\textit{Source reference.} Original lemma 079W: sources/stacks-chow/tags.txt, source lines 9286--9286. Cross-chapter standard prerequisite located in the project tag list staged in the local snapshot; invoked at proof line 3083 of the counted proof. It is retained only as a navigation pointer; every load-bearing step is inside the delivered span.\par\medskip
\ifdefined\hypertarget\hypertarget{sourcepointer-8}{}\fi
\noindent\textit{Source reference.} Original lemma 08EU: sources/stacks-chow/tags.txt, source lines 10685--10685. Cross-chapter standard prerequisite located in the project tag list staged in the local snapshot; invoked at proof line 3086 of the counted proof. It is retained only as a navigation pointer; every load-bearing step is inside the delivered span.\par\medskip
\ifdefined\hypertarget\hypertarget{sourcepointer-9}{}\fi
\noindent\textit{Source reference.} Original lemma 0658: sources/stacks-chow/tags.txt, source lines 7892--7892. Cross-chapter standard prerequisite located in the project tag list staged in the local snapshot; invoked at proof line 3093 of the counted proof. It is retained only as a navigation pointer; every load-bearing step is inside the delivered span.\par\medskip
\ifdefined\hypertarget\hypertarget{sourcepointer-10}{}\fi
\noindent\textit{Source reference.} Original lemma 08DJ: sources/stacks-chow/tags.txt, source lines 10639--10639. Cross-chapter standard prerequisite located in the project tag list staged in the local snapshot; invoked at proof line 3123 of the counted proof. It is retained only as a navigation pointer; every load-bearing step is inside the delivered span.\par\medskip
\ifdefined\hypertarget\hypertarget{sourcepointer-11}{}\fi
\noindent\textit{Source reference.} Original lemma 0A6H: sources/stacks-chow/tags.txt, source lines 12845--12845. Cross-chapter standard prerequisite located in the project tag list staged in the local snapshot; invoked at proof line 3143 of the counted proof. It is retained only as a navigation pointer; every load-bearing step is inside the delivered span.\par\medskip
\ifdefined\hypertarget\hypertarget{sourcepointer-12}{}\fi
\noindent\textit{Source reference.} Original theorem 09IS: sources/stacks-chow/tags.txt, source lines 12050--12050. Cross-chapter standard prerequisite located in the project tag list staged in the local snapshot; invoked at proof line 3154 of the counted proof. It is retained only as a navigation pointer; every load-bearing step is inside the delivered span.\par\medskip
\begin{lemma}
\label{lemma-properties-relative-dualizing}
Let $Y$ be a quasi-compact and quasi-separated scheme.
Let $f : X \to Y$ be a morphism of schemes which is
proper, flat, and of finite presentation with
relative dualizing complex $\omega_{X/Y}^\bullet$
(Remark \ref{remark-relative-dualizing-complex}).
Then
\begin{enumerate}
\item $\omega_{X/Y}^\bullet$ is a $Y$-perfect object of $D(\mathcal{O}_X)$,
\item $Rf_*\omega_{X/Y}^\bullet$ has vanishing cohomology sheaves
in positive degrees,
\item $\mathcal{O}_X \to
R\SheafHom_{\mathcal{O}_X}(\omega_{X/Y}^\bullet, \omega_{X/Y}^\bullet)$
is an isomorphism.
\end{enumerate}
\end{lemma}

\begin{proof}
In view of the fact that formation of $\omega_{X/Y}^\bullet$ commutes
with base change (see Remark \ref{remark-relative-dualizing-complex}),
we may and do assume that $Y$ is affine. For a perfect object $E$ of
$D(\mathcal{O}_X)$ we have
\begin{align*}
Rf_*(E \otimes_{\mathcal{O}_X}^\mathbf{L} \omega_{X/Y}^\bullet)
& =
Rf_*R\SheafHom_{\mathcal{O}_X}(E^\vee, \omega_{X/Y}^\bullet) \\
& =
R\SheafHom_{\mathcal{O}_Y}(Rf_*E^\vee, \mathcal{O}_Y) \\
& =
(Rf_*E^\vee)^\vee
\end{align*}
For the first equality, see
Cohomology, Lemma \ref{cohomology-lemma-dual-perfect-complex}.
For the second equality, see Lemma \ref{lemma-iso-on-RSheafHom},
Remark \ref{remark-iso-on-RSheafHom}, and 
Derived Categories of Schemes, Lemma
\ref{perfect-lemma-flat-proper-perfect-direct-image-general}.
The third equality is the definition of the dual. In particular
these references also show that the outcome is a perfect object
of $D(\mathcal{O}_Y)$. We conclude that $\omega_{X/Y}^\bullet$
is $Y$-perfect by More on Morphisms, Lemma
\ref{more-morphisms-lemma-characterize-relatively-perfect}.
This proves (1).

\medskip\noindent
Let $M$ be an object of $D_\QCoh(\mathcal{O}_Y)$. Then
\begin{align*}
\Hom_Y(M, Rf_*\omega_{X/Y}^\bullet) & =
\Hom_X(Lf^*M, \omega_{X/Y}^\bullet) \\
& =
\Hom_Y(Rf_*Lf^*M, \mathcal{O}_Y) \\
& =
\Hom_Y(M \otimes_{\mathcal{O}_Y}^\mathbf{L} Rf_*\mathcal{O}_X, \mathcal{O}_Y)
\end{align*}
The first equality holds by Cohomology, Lemma
\ref{cohomology-lemma-adjoint}.
The second equality by construction of $a$.
The third equality by Derived Categories of Schemes, Lemma
\ref{perfect-lemma-cohomology-base-change}.
Recall $Rf_*\mathcal{O}_X$ is perfect of tor amplitude in $[0, N]$
for some $N$, see
Derived Categories of Schemes, Lemma
\ref{perfect-lemma-flat-proper-perfect-direct-image-general}.
Thus we can represent $Rf_*\mathcal{O}_X$ by a complex of
finite projective modules sitting in degrees $[0, N]$
(using More on Algebra, Lemma \ref{more-algebra-lemma-perfect}
and the fact that $Y$ is affine).
Hence if $M = \mathcal{O}_Y[-i]$ for some $i > 0$, then the last
group is zero. Since $Y$ is affine we conclude that
$H^i(Rf_*\omega_{X/Y}^\bullet) = 0$ for $i > 0$.
This proves (2).

\medskip\noindent
Let $E$ be a perfect object of $D_\QCoh(\mathcal{O}_X)$. Then
we have
\begin{align*}
\Hom_X(E, R\SheafHom_{\mathcal{O}_X}(\omega_{X/Y}^\bullet, \omega_{X/Y}^\bullet)
& =
\Hom_X(E \otimes_{\mathcal{O}_X}^\mathbf{L} \omega_{X/Y}^\bullet,
\omega_{X/Y}^\bullet) \\
& =
\Hom_Y(Rf_*(E \otimes_{\mathcal{O}_X}^\mathbf{L} \omega_{X/Y}^\bullet),
\mathcal{O}_Y) \\
& =
\Hom_Y(Rf_*(R\SheafHom_{\mathcal{O}_X}(E^\vee, \omega_{X/Y}^\bullet)),
\mathcal{O}_Y) \\
& =
\Hom_Y(R\SheafHom_{\mathcal{O}_Y}(Rf_*E^\vee, \mathcal{O}_Y),
\mathcal{O}_Y) \\
& =
R\Gamma(Y, Rf_*E^\vee) \\
& =
\Hom_X(E, \mathcal{O}_X)
\end{align*}
The first equality holds by Cohomology, Lemma
\ref{cohomology-lemma-internal-hom}.
The second equality is the definition of $\omega_{X/Y}^\bullet$.
The third equality comes from the construction of the dual perfect
complex $E^\vee$, see Cohomology, Lemma
\ref{cohomology-lemma-dual-perfect-complex}.
The fourth equality follows from the equality
$Rf_*R\SheafHom_{\mathcal{O}_X}(E^\vee, \omega_{X/Y}^\bullet) =
R\SheafHom_{\mathcal{O}_Y}(Rf_*E^\vee, \mathcal{O}_Y)$
shown in the first paragraph of the proof.
The fifth equality holds by double duality for perfect complexes
(Cohomology, Lemma
\ref{cohomology-lemma-dual-perfect-complex})
and the fact that $Rf_*E$ is perfect by
Derived Categories of Schemes, Lemma
\ref{perfect-lemma-flat-proper-perfect-direct-image-general}.
The last equality is Leray for $f$.
This string of equalities essentially shows (3)
holds by the Yoneda lemma. Namely, the object
$R\SheafHom(\omega_{X/Y}^\bullet, \omega_{X/Y}^\bullet)$
is in $D_\QCoh(\mathcal{O}_X)$ by Derived Categories of Schemes, Lemma
\ref{perfect-lemma-quasi-coherence-internal-hom}.
Taking $E = \mathcal{O}_X$ in the above we get a map
$\alpha : \mathcal{O}_X \to
R\SheafHom_{\mathcal{O}_X}(\omega_{X/Y}^\bullet, \omega_{X/Y}^\bullet)$
corresponding to
$\text{id}_{\mathcal{O}_X} \in \Hom_X(\mathcal{O}_X, \mathcal{O}_X)$.
Since all the isomorphisms above are functorial in $E$ we
see that the cone on $\alpha$ is an object $C$ of $D_\QCoh(\mathcal{O}_X)$
such that $\Hom(E, C) = 0$ for all perfect $E$.
Since the perfect objects generate
(Derived Categories of Schemes, Theorem
\ref{perfect-theorem-bondal-van-den-Bergh})
we conclude that $\alpha$ is an isomorphism.
\end{proof}

\end{document}
